To avoid potential confusion we remark that throughout we will call any matrix $A\in \FF^{n_1} \otimes \FF^{n_2}$ (that is not necessarily square) \emph{diagonal} if its support, which is defined as $\supp(A) := \{(i,j) : A_{i,j} \neq 0\}$, is a subset of $\{(i,i): 1\leq i\leq \min\{n_1, n_2\}\}$.
%

%
%


\subsection{Min-rank of matrix subspaces and its properties}\label{subsec:minrank} %
%
%

%
\newcommand{\StrExC}{T_{\textnormal{NA},c,n}}
%
%
%
%
%
%
%

%
%
%

%

%

%

%
The min-rank of a matrix subspace is simply the smallest rank of any non-zero matrix in the space. 


\begin{definition}
    %
    We define the \emph{min-rank} of a matrix subspace $\mathcal{A} \subseteq \FF^{n_1} \otimes \FF^{n_2}$ as
    %
    \[
    \minrnk(\mathcal{A}):=\min\{\rank(A):A\in\mathcal{A}\setminus\{0\}\}.
    \]
    For any collection of matrices $A_1,
    \ldots,A_c \in \FF^{n_1} \otimes \FF^{n_2}$ we define \[
    \minrnk(A_1,\dots,A_c):=\minrnk(\linspan\{A_1,\dots,A_c\}).
    \]
\end{definition}



%
%
%

%
%
%
%

We will use the following straightforward equivalence. % throughout this section.
%Given matrices $A_1, \ldots, A_c \in \FF^{n_1} \otimes \FF^{n_2}$ we call a linear combination $\alpha_1 A_1 + \cdots + \alpha_c A_c$ with coefficients $\alpha_i \in \FF$ \emph{nontrivial} if at least one of the coefficients $\alpha_i$ is nonzero.

\begin{lemma}\label{lem:minrank-nontrivial-lin-comb}
    Let $A_1,\dots,A_c \in \FF^{n_1} \otimes \FF^{n_2}$ be matrices. %
    For every $r\geq1$, the following two statements are equivalent:
    \begin{enumerate}[\upshape(a)]
        \item The matrices $A_1,\dots,A_c$ are linearly independent and $\minrnk(A_1,\dots,A_c)\geq r$.
        \item For every $u \in \FF^c \setminus \{0\}$ we have $\rank(u_1 A_1 + \cdots + u_c A_c) \geq r$.
    \end{enumerate} 
\end{lemma}
